\section{The problem: a regulatory decision loses meaning as it becomes software}
\label{sec:problem}

A private bank must turn regulatory requirements into client classifications,
product permissions, disclosures, monitoring, records and staff actions. Those
requirements rarely arrive as a complete algorithm. A circular may change one
part of an existing regime, depend on definitions in a code, distinguish product
providers from distributors, or explain an exception in an annex. Software must
nevertheless decide what to do for a particular client, product and event.

The sponsor's reported problem is the error-prone manual reading and updating of
bank systems. This proposal treats that account as the starting business need;
it does not assume a measured incident rate. The process below is a plausible
operating model to validate with the bank. A compliance officer reads a circular,
writes an interpretation, and discusses it with product and operations staff. A
business analyst turns the interpretation into requirements. Engineers implement
configuration or code, and testers construct examples from those requirements.
Each participant can perform competently while a condition disappears between
documents. When the circular changes, the bank must reconstruct the chain.

\begin{figure}[htbp]
\centering
\begin{tikzpicture}[node distance=5mm, every node/.style={font=\small},
  box/.style={draw=teal,rounded corners,fill=pale,align=center,text width=25mm,minimum height=15mm},
  arr/.style={-{Latex[length=2mm]},thick,draw=navy}]
\node[box] (a) {Circulars, annexes, codes};
\node[box,right=of a] (b) {Compliance interpretation};
\node[box,right=of b] (c) {Requirements and tickets};
\node[box,right=of c] (d) {System rules and workflow};
\node[box,right=of d] (e) {Client decisions and records};
\draw[arr] (a)--(b); \draw[arr] (b)--(c); \draw[arr] (c)--(d); \draw[arr] (d)--(e);
\end{tikzpicture}
\caption{The proposed discovery exercise will map the bank's actual process onto
these transitions. LegalMath preserves the source, interpretation and tests across
them, so a reviewer can follow a decision back to its reason.}
\label{fig:manual}
\end{figure}

The product opportunity is to make that chain inspectable and reusable. A
reviewer should be able to select a rule and see the passage that supports it,
the interpretation adopted, the facts it needs, the situations in which it changes
an outcome, and the version actually used. A subsequent amendment should identify
the affected controls and test cases without relying on the memory of the
original implementation team.

\subsection{The regulatory perimeter is part of the specification}

The supplied Hong Kong Securities and Futures Commission (SFC) product-authorization
page is a useful discovery source, but it
is not a complete inventory of a private bank's distribution obligations. The
SPI guidance examined here is a joint circular from the SFC and Hong Kong Monetary
Authority (HKMA) to intermediaries. The
tokenisation circular distinguishes product-provider responsibilities from
distributor responsibilities and imports other product and intermediary rules.
The marketing circular refers to the Securities and Futures Ordinance, the Code
of Conduct and advertising guidance. These dependencies determine the meaning of
the extracted clauses \citep{sfcindex,sfcspi,sfctoken2026,sfcmarketing}.

Discovery therefore starts with the bank's legal entities and activities. For a
specified business, the inventory should record whether the institution acts as
distributor, product provider, adviser or another regulated actor, and which
client and product classes are involved. Being a Hong Kong office of a US banking
group does not by itself settle any of these questions. The prototype will record
the selected Hong Kong perimeter and leave other group and jurisdictional
requirements as separately identified dependencies. A broader rollout requires
the bank's authoritative inventory.

Document classification also matters. The September 2025 joint circular announces
a concurrent thematic review of distribution practices and describes what the
review will examine. It should enter regulatory-change tracking; its description
of the review is not, by itself, a new numerical client-eligibility test
\citep{sfcreview2025}. A translator that treats every sentence as a binding
conditional will manufacture rules. Conversely, a translator that extracts only
sentences containing a particular modal verb can miss definitions, incorporated
requirements and material explanatory guidance.

LegalMath should preserve a passage's role: definition, requirement, permission,
exception, supervisory expectation, explanation or announcement. The original
wording and its authority remain alongside the adopted implementation. A bank may
choose an additional operational restriction, but that choice should be recorded
as bank policy with an owner and rationale. It must not be silently attributed to
the SFC.

\subsection{A running case: sophisticated professional investors}
\label{sec:spi}

The joint circular of 28 July 2023 supplies an unusually useful test case because
its apparent simplicity conceals several different kinds of rule. Annex 1
defines the qualifying criteria, eligible investment transactions, information
and consent requirements, and continuing controls. Annex 2 clarifies difficult
implementation cases. The examples below are synthetic and illustrate the cited
provisions; they do not complete an SPI assessment or authorize a transaction
\citep{sfcspi,sfcspiannex1,sfcspiannex2}.

\caseheading{A one-word error changes the financial condition}
Annex 1, paragraph 3.1 provides two alternative financial routes: a portfolio of
at least HK\$40 million, or net assets excluding the primary residence of at least
HK\$80 million, including the stated foreign-currency equivalents. Consider a
client with an eligible portfolio valued at HK\$35 million and qualifying net
assets of HK\$90 million. The second route meets this financial condition.
Replacing the alternative with a conjunction rejects the client incorrectly.
Replacing either inclusive comparison with a strict comparison rejects a client
exactly at that threshold. Both defects are easy to miss when test data use
wealth amounts far from the boundary \citep[Annex 1, para.~3.1]{sfcspiannex1}.

The numbers also require definitions. The value of assets held at the bank is
not automatically the qualifying portfolio, and total wealth is not automatically
net assets excluding the primary residence. Paragraph 3.3 specifies the net-asset
calculation. In a simple synthetic case with HK\$120 million of assets, including
a HK\$30 million primary residence, and HK\$15 million of liabilities, that
calculation gives HK\$75 million: $120-15-30$. Omitting the residence gives
HK\$105 million and changes the result of the net-asset route. The portfolio
route must still be assessed independently. The bank must also establish how the
source definitions map to its valuation and ownership records
\citep[Annex 1, paras.~3.2--3.4]{sfcspiannex1}.

\caseheading{Ownership and missing evidence can reverse the result}
For a joint portfolio with persons other than the client's associate, paragraph
3.2(c) refers to the client's share specified in a written agreement, or an equal
share in the absence of that agreement. Suppose a two-holder portfolio is worth
HK\$60 million and the written allocation gives the client 20 percent. The
relevant share for this illustration is HK\$12 million, not the whole account
and not an automatically assumed half. The rule for a joint account with an
associate is separately stated. A generic ``joint-account multiplier'' loses
this distinction \citep[Annex 1, para.~3.2]{sfcspiannex1}.

An unavailable document is also different from an established absence of an
agreement. The bank must decide what evidence establishes each fact. If a feed
returns a null portfolio value, replacing null with zero creates a substantive
conclusion from a data failure. If qualifying net assets are known to be
HK\$90 million, the financial condition can still be established by that route.
If they are HK\$70 million, the unknown portfolio value can change the answer;
the financial condition remains unresolved. The product must express that
distinction directly.

\caseheading{Wealth does not establish suitability for streamlined treatment}
Paragraph 1 combines professional-investor status, the financial criteria,
knowledge or experience, and investment objectives at the relevant date.
Paragraph 4.1 requires reasonable satisfaction concerning sophistication and
sets out alternative supporting criteria. Paragraph 5.1 excludes conservative
clients from SPI treatment for this guidance. A wealthy client whose objective
is capital preservation therefore demonstrates why a field named
\texttt{wealth\_qualified} must not automatically populate
\texttt{streamlined\_approach\_available}
\citep[Annex 1, paras.~1, 4--5]{sfcspiannex1}.

Some criteria are countable; others require assessment. A degree, a qualification,
relevant work experience and prior transactions are distinct routes. A translator
should not invent a score that trades off wealth against sophistication or
investment objectives. The assessment and its supporting evidence need an owner,
date and review record. Paragraph 2 also supplies a separate route for a qualifying
investment-holding corporation wholly owned by SPIs. It cannot be represented by
treating every client identifier as an individual.

\caseheading{Experience and consent belong to a product category}
Suppose the client relies on the transaction-experience route and has completed
five bond transactions in the relevant three-year period. That history does not
automatically establish the same route for an unrelated product category.
Paragraph 7.3 connects category eligibility to the alternative qualification and
experience routes; paragraph 7.2 requires the client's choice of categories to be
documented and category information to be provided. A global client flag can
therefore remain true while a proposed category is outside the selected scope
\citep[Annex 1, paras.~4.1, 7.1--7.3]{sfcspiannex1}.

The identity of the person matters too. Annex 2, FAQ 2 explains the assessment
where a client acts through a power of attorney. The attorney's sophistication
cannot simply replace the client's qualifying attributes. In software terms,
the customer, account holder, representative, beneficial owner and person giving
an instruction are different roles. Collapsing them into one record can make
every subsequent logical step internally consistent and still apply it to the
wrong person \citep[Annex 2, FAQ 2]{sfcspiannex2}.

\caseheading{Exposure monitoring is an event rule}
Paragraph 8.3 permits two monitoring approaches. One concerns gross exposure on
execution of transactions under the streamlined approach. The other uses
designated accounts and checks the specified condition after a top-up or deposit
of new funds. Annex 2, FAQ 6 explains how these approaches operate, including
transferred exposures and restrictions on additional funds when a designated
account exceeds the threshold. FAQ 7 addresses leverage
\citep[Annex 1, para.~8.3; Annex 2, FAQs 6--7]{sfcspiannex1,sfcspiannex2}.

A single rule that rejects every transaction whenever an account's current value
exceeds a threshold loses the distinction between these approaches. In the
designated-account example, the FAQ allows continued execution above the
threshold subject to its conditions restricting new funds. The monitoring mode,
movement of funds, gross exposure and prior events all matter. An appreciated
position, a leveraged purchase, a cash withdrawal and a transferred non-cash
position must not be reduced to the same change in account market value. This is
a strong reason to model event histories as well as decision tables.

\caseheading{Consent and continuing obligations outlive the initial decision}
Paragraph 13 calls for prior written agreement and an explanation of the ability
to withdraw from, or change, the arrangement. Paragraph 14 requires annual review
and reminders. Assessment and agreement are therefore not permanent facts. A
transaction considered before receipt of a withdrawal instruction and one
considered afterwards require different state snapshots. A review reminder being
queued is different from being sent; being sent is different from the review
being completed \citep[Annex 1, paras.~12--14]{sfcspiannex1}.

The streamlined approach also retains specific information and conduct
requirements. The solicited or recommended transaction provisions and the
unsolicited complex-product provisions have different contexts. A broad
``skip suitability and disclosure'' switch would conceal these remaining duties
\citep[Annex 1, paras.~9--11]{sfcspiannex1}. The prototype should return the
relevant decision and the associated tasks, explanations and evidence requests,
rather than a single apparently comprehensive compliance flag.

\subsection{Amendment, actor and conditional language: tokenised products}

The SFC's 20 April 2026 tokenisation circular replaces the November 2023 circular;
the earlier public entry is marked accordingly. A source collector that keeps
only a URL and the most recently extracted text cannot reliably reconstruct
which requirements an older decision used. A collector that retains both but
fails to record replacement can apply inconsistent generations together
\citep{sfctoken2023,sfctoken2026}.

The 2026 circular also provides useful clause-level challenges. Paragraph 10
keeps responsibility for the tokenisation arrangement and ownership records
with product providers despite outsourcing. Paragraph 18 addresses distributors,
including providers distributing their own products. These are actor-specific
requirements: the word ``bank'' is not a sufficient applicability condition.
Paragraph 13 conditions the use of public-permissionless networks on additional
proper controls; translating it into an unconditional technology ban changes the
statement \citep[paras.~10, 13, 18]{sfctoken2026}.

Paragraphs 14--16 illustrate several different triggers. Confirmation and
demonstration upon request are distinguishable actions. Third-party verification
and a legal opinion are conditioned on an SFC request in the cited paragraphs.
Removing that trigger creates a different obligation; removing the obligation
creates the opposite error. A request must be represented as an event, with its
subject and responsible actor, and any deadline must come from a relevant source
or separately identified operational policy. The product must not invent a
generic deadline \citep[paras.~14--16]{sfctoken2026}.

Prior consultation and prior approval are also different. Paragraphs 20--21 and
footnote 8 distinguish new products, tokenisation of existing products, material
changes, and circumstances concerning tokenised share classes. Turning
``consultation required'' into ``approval obtained'' skips an action and invents
a result. The proposed workbench will expose such verbs and their dependencies
as separate obligations and evidence fields \citep[paras.~20--21 and n.~8]{sfctoken2026}.

\subsection{Marketing rules require exceptions and judgment}

The October 2023 distributor circular examines additional returns and incentives,
redemption restrictions and marketing of fund-related services. Its treatment of
guaranteed-return arrangements depends on the arrangement and the relevant legal
classification. A potential structured-product classification should produce a
supported assessment task; a language model should not resolve it from a slogan
alone \citep[paras.~5--11]{sfcmarketing}.

The circular's discussion of gifts expressly distinguishes discounts of fees or
charges. Losing that exception makes an apparently conservative implementation
incorrectly broad. Its redemption discussion distinguishes administrative
procedures or cut-off times from reducing a fund's dealing frequency. A daily
dealing fund offered with weekly redemption presents a different issue from a
different daily processing cut-off. These are excellent paired test cases because
they differ in the fact that determines the treatment
\citep[paras.~9--14]{sfcmarketing}.

The advertising discussion addresses misleading impressions and comparisons with
deposits. Keyword detection can nominate material for review, but it cannot
establish the whole classification: meaning depends on the presentation and the
service described. The system should retain the advertisement, the relevant
source passages and the reviewer's conclusion. That is a useful result even when
there is no defensible fully automated predicate
\citep[paras.~15--17]{sfcmarketing}.

\subsection{What makes the current work expensive and hard to verify}

The common failure is loss of context. Thresholds detach from definitions;
exceptions detach from their parent rules; duties detach from actors and events;
and historical decisions detach from the source version that justified them.
Document search helps locate text, but it does not preserve those relationships
through implementation. More fluent summaries can make an omitted condition
harder to notice.

Review then becomes repetitive. Compliance officers must infer what an engineer
understood, engineers must infer what a compliance memo leaves implicit, and
testers must construct enough cases to expose mismatches. If tests are written
only from the implementation ticket, a mistranslation shared by the ticket and
code can pass every test. After an amendment, teams may retest too little because
dependencies are hidden, or retest too much because they cannot show what is
unaffected.

The proposed product makes each material interpretation a maintained part of the
system. Its value should be measured in fewer material defects, more complete
explanations, and less repeated review effort. Those benefits are hypotheses for
the prototype to test. The evidence needed to test them includes independently
reviewed cases, disagreements that remain visible, and the time spent resolving
them; a high extraction score alone cannot establish the business case.
